\documentclass[11pt]{article}
\usepackage[margin=0.9in]{geometry}
\usepackage{amsmath,amssymb,amsthm}
\usepackage[T1]{fontenc}
\usepackage{lmodern}
\usepackage[hidelinks]{hyperref}
\newtheorem{theorem}{Theorem}
\newtheorem{lemma}{Lemma}
\newcommand{\Rea}{\operatorname{Re}}
\newcommand{\Ima}{\operatorname{Im}}
\newcommand{\cross}[2]{[#1,#2]}
\title{The outer-pedal area centroid in even-period elliptic billiards}
\author{Alec Kriebel\\Independent researcher\\
\small\href{https://orcid.org/0009-0001-9320-500X}{ORCID: 0009-0001-9320-500X}}
\date{13 October 2026 \quad (version 1.0, unrefereed)}
\begin{document}
\maketitle
\begin{abstract}
Reznik, Garcia and Koiller observed that the area centroid of the outer-side
pedal polygon of an even-period elliptic billiard is independent of the orbit
phase (invariant k307). We prove this observation on the natural domain of the
signed-area centroid and give its explicit value. The proof uses antipodal
pairing, support coordinates and telescoping edge moments. For simple even
orbits the centroid exists for every fixed pedal point. For even star orbits
it can be undefined on a fixed circle, which we determine exactly and realize
in a noncircular eight-period example. The note includes a reproducible
verification package; the written identities give the all-period proof.
\end{abstract}

\section{The statement and its domain}
Let the billiard boundary and its fixed caustic be
\[
 E:\frac{x^2}{a^2}+\frac{y^2}{b^2}=1,\qquad
 E_c:\frac{x^2}{a^2-\lambda}+\frac{y^2}{b^2-\lambda}=1,
 \qquad a>b>0,\quad 0<\lambda<b^2.
\]
Assume this pair supports a Poncelet billiard family of least even period
$N\geq4$. Primitive star trajectories are allowed. The vertices $P_i$ are
ordered along the trajectory, with indices modulo $N$. The outer polygon has
as its side lines the tangents to $E$ at the $P_i$. For a fixed point $M$, let
$Q_i$ be its perpendicular foot on that full tangent line.
For any ordered polygon $q_i$, identify the plane with $\mathbb C$ and set
\begin{equation}\label{signed}
 S(q)=\frac12\sum_i\cross{q_i}{q_{i+1}},\qquad
 \mathcal T(q)=\sum_i\cross{q_i}{q_{i+1}}(q_i+q_{i+1}),\qquad
 C(q)=\frac{\mathcal T(q)}{6S(q)},
\end{equation}
where $\cross{p}{q}=\Ima(\overline p q)$, and $C$ is defined only if $S\ne0$.
These are signed moments, including multiplicity for self intersections.
They agree with the area-centroid convention of \cite{RGK}; they are neither
the arithmetic vertex mean nor an unsigned union-of-regions centroid.

Put
\begin{align*}
 H&=(a^2+b^2)/2,& e&=(a^2-b^2)/2,\\
 d&=1-\lambda(a^{-2}+b^{-2}),&F&=Hd+2\lambda,\\
 v&=\frac{-4e\lambda d}{F^2-e^2d^2},&z&=M_x+iM_y,\\
 R^2&=|z|^2,&G&=2F+dR^2.
\end{align*}
Write $S(M)=S(Q(M))$ for the chosen member of the family.
\begin{theorem}\label{main}
Orient the billiard counterclockwise around its elliptic caustic. Then
$F>0$, $F^2-e^2d^2>0$, $S(0)>0$, and
\begin{equation}\label{area}
 S(M)=S(0)\frac{2F+d|z|^2}{2F}.
\end{equation}
Whenever $G\ne0$, the signed area centroid is
\begin{equation}\label{centroid}
 C(Q(M))=\frac z2-
 \frac{z(F+2Hd+evd)+\overline z(2Fv+3ed+vd|z|^2/2)}
 {3(2F+d|z|^2)}.
\end{equation}
Thus it is independent of the orbit phase wherever defined.
If $d\geq0$, it is defined for every $M$. If $d<0$, its exact undefined locus
is the phase-independent circle $|M|^2=-2F/d$. Least period four is equivalent
to $d=0$, when the centroid is $M/3$. Every simple winding-one even family
has $d\geq0$.
\end{theorem}

The theorem proves the k307 observation in Table 4 and Section 3.4 of
\cite{RGK}, with the denominator qualification already inherent in its signed
centroid definition. An unqualified assertion that a centroid exists for
every $M$ in every even star family would be false. Equation \eqref{area}
does not assert that $S(0)$ is phase independent.
Reversal changes both signed moments by a sign and leaves the centroid
unchanged. Repeated traversals of an even least-period orbit give the same
centroid; repeating an odd primitive orbit does not create the antipodal
hypothesis. Hyperbolic and degenerate caustics, and the circular boundary
case, are outside the theorem.

The originating observation is due to Reznik--Garcia--Koiller.
Bialy--Tabachnikov \cite[Theorems 4.1 and 4.3]{BT} prove the neighboring
vertex-centroid invariant k306 and recall the classical even-period symmetry.
Schwartz--Tabachnikov \cite{ST} study signed centroid loci of Poncelet
polygons between conics. We do not claim novelty for those methods, the
symmetry, or the cyclic lemma below. Two independent bounded primary-source
audits completed on 13 October 2026 found no earlier proof covering the
present full statement; this is not a certification of historical firstness.

\section{Support coordinates and antipodal pairing}
Let $T$ be the counterclockwise Poncelet map on $E$ with the specified tangent
choice, and $A(P)=-P$ its half-turn. The porism gives $T^N=\mathrm{id}$,
and $A$ commutes with $T$. The finite group generated by them consists of
orientation-preserving circle homeomorphisms and is cyclic. Indeed a finite
order orientation-preserving circle homeomorphism fixing a point is the
identity: after cutting there, a moved point has strictly monotone iterates.
The group therefore acts faithfully on a finite orbit, preserving its circular
order, and embeds into its cyclic permutation group. Uniqueness of the
involution gives $T^{N/2}=A$. Consequently
\begin{equation}\label{antipodes}
 P_{i+N/2}=-P_i.
\end{equation}
Let $n_i=(\cos\psi_i,\sin\psi_i)$ be the outward unit normal and $h_i>0$
the support number. Then
\[
 h_i^2=H+e\cos2\psi_i,\qquad
 n_{i+N/2}=-n_i,\qquad h_{i+N/2}=h_i.
\]
The caustic strictly contains the origin. Each oriented tangent chord has its
next endpoint less than a half-turn ahead on $E$. The increasing Gauss map
commutes with the half-turn, so successive lifted normal angles satisfy
$0<\Delta_i=\psi_{i+1}-\psi_i<\pi$. In particular the outer tangent
intersections are finite, and
\begin{equation}\label{positive}
 S(0)=\frac12\sum_i h_i h_{i+1}\sin\Delta_i>0.
\end{equation}
This uses trajectory order and applies to stars too.

For consecutive eccentric angles $u,w$, put $\tau=(u+w)/2$ and
$q=(w-u)/2\in(0,\pi/2)$. Their chord is
$(\cos\tau/a)x+(\sin\tau/b)y=\cos q$. Tangency to $E_c$ gives
$\sin^2q=\lambda(\cos^2\tau/a^2+\sin^2\tau/b^2)$, equivalently
\[
 d=[1+\lambda(a^{-2}-b^{-2})]\cos u\cos w
  +[1-\lambda(a^{-2}-b^{-2})]\sin u\sin w.
\]
Since the normal is $h_i(\cos u/a,\sin u/b)$, this yields
\begin{equation}\label{support}
 d h_i h_{i+1}=F\cos\Delta_i+ed\cos(\psi_i+\psi_{i+1}).
\end{equation}
All denominators in the theorem are controlled by
\[
 F+ed=a^2-\lambda(a^2/b^2-1)>b^2>0,\qquad
 F-ed=b^2+\lambda(1-b^2/a^2)>0.
\]
In particular $F>0$ and $F^2-e^2d^2>0$.

If $N=4$, apply \eqref{support} to $(n_0,n_1)$ and $(n_1,n_2)$.
Their right sides negate because $n_2=-n_0$, whereas both left sides are
$d h_0h_1$. Hence $2d h_0h_1=0$ and $d=0$.
Conversely $d=0$ makes $\cos\Delta_i=0$, so each normal advances $\pi/2$
and the least period is four. Its outer polygon is a rectangle. In an
orthonormal side frame its pedal vertices are
\[
 (A,M_y),\ (M_x,B),\ (-A,M_y),\ (M_x,-B),\qquad A,B>0.
\]
Direct shoelace gives $S=2AB$ and $\mathcal T=4ABM$, proving both formulas
and $C=M/3$. For the remaining proof assume $d\ne0$ and set $D=F/d$.
Then $D^2-e^2>0$ and \eqref{support} reads
\begin{equation}\label{Dsupport}
 h_i h_{i+1}=D\cos\Delta_i+e\cos(\psi_i+\psi_{i+1}).
\end{equation}

\section{A doubled-normal telescope}
Regard $n_i$ as a unit complex number and put $Z_i=n_i^2$.
Squaring \eqref{Dsupport} and substituting the support squares gives
\[
 (D^2-e^2)\cos^2\Delta_i
 =H^2-e^2+e(H-D)(\cos2\psi_i+\cos2\psi_{i+1}).
\]
Thus, with $u=2(H^2-e^2)/(D^2-e^2)-1$,
\begin{equation}\label{cosrelation}
 \Rea(\overline Z_iZ_{i+1})=u+v(\Rea Z_i+\Rea Z_{i+1}),\qquad
 v=\frac{2e(H-D)}{D^2-e^2}=\frac{-4e\lambda d}{F^2-e^2d^2}.
\end{equation}
\begin{lemma}\label{cyclic}
For ordered unit complex vertices satisfying \eqref{cosrelation},
$\mathcal T(Z)=2vS(Z)$. No simplicity or nonzero-area assumption is needed.
\end{lemma}
\begin{proof}
For unit $U,V$ the division-free identity
\[
 \cross{U}{V}(U+V)=-i(1+\Rea(\overline U V))(V-U)
\]
follows by multiplication. The constant part telescopes. Write
$Z_i=X_i+iY_i$. The real part of
$\sum_i(X_i+X_{i+1})(Z_{i+1}-Z_i)$ telescopes as
$\sum_i(X_{i+1}^2-X_i^2)=0$. Its imaginary part is
$\sum_i\cross{Z_i}{Z_{i+1}}$, because
\[
 (X_i+X_{i+1})(Y_{i+1}-Y_i)
 =\cross{Z_i}{Z_{i+1}}+X_{i+1}Y_{i+1}-X_iY_i.
\]
Multiplication by $-iv$ proves the assertion.
\end{proof}
Write $S_Z=S(Z)$ and $\mathcal T_Z=\mathcal T(Z)$.
The origin-area identity is
\begin{equation}\label{origin}
 S(0)=\frac D2 S_Z.
\end{equation}
Indeed multiply \eqref{Dsupport} by $\sin\Delta_i$ and use
$\cos(\psi_i+\psi_{i+1})\sin\Delta_i
=(\sin2\psi_{i+1}-\sin2\psi_i)/2$; the latter terms telescope.
By \eqref{positive}, $S_Z\ne0$ when $d\ne0$.

\section{The pedal area and full first moment}
Projection onto $n_i\cdot x=h_i$ gives
$Q_i=M+(h_i-M\cdot n_i)n_i$. Put $t=\overline z/2$ and $w_i=h_i n_i$.
In complex notation,
\begin{equation}\label{projection}
 Q_i=z/2+V_i,\qquad V_i=w_i-tZ_i.
\end{equation}
Opposite indices negate $w_i$ and fix $Z_i$; hence terms odd in the $w$'s
cancel in cyclic sums. Area is translation invariant, so
\begin{equation}\label{pedalarea}
 S(M)=S(0)+|t|^2S_Z=(D/2+|z|^2/4)S_Z,
\end{equation}
which proves \eqref{area}.

We give the complete cubic reduction to make the essential calculation
checkable. In $\mathcal T(V)$, the pure-$w$ and quadratic-in-$t$ terms
cancel by antipodal pairing. The cubic term is $-|t|^2t\mathcal T_Z$.
For one edge put $x=n_i$, $y=n_{i+1}$, $h=h_i$, $k=h_{i+1}$, $Z=x^2$,
$W=y^2$. Average the edge moment with its antipodal edge moment.
In $2i$ times that average the coefficients of
$t$ and $\overline t$ are, respectively,
\begin{align*}
 c&=k^2x^2-h^2y^2+2hk(x^3/y-y^3/x),\\
 b&=h^2x^2/y^2-k^2y^2/x^2+hk(x/y-y/x).
\end{align*}
Substitute
\[
 h^2=H+\tfrac e2(x^2+x^{-2}),\quad
 k^2=H+\tfrac e2(y^2+y^{-2}),\quad
 hk=\tfrac D2(y/x+x/y)+\tfrac e2(xy+(xy)^{-1}).
\]
With $r=Z/W-W/Z$ and $s=Z^2/W-W^2/Z$, multiplication gives the exact
edge identities
\begin{align*}
 c&=Ds+\tfrac{3e}{2}r+(D+H)(Z-W)+e(Z^2-W^2),\\
 b&=(D/2+H)r+\tfrac e2s+\tfrac e2(Z-W)+e(W^{-1}-Z^{-1}).
\end{align*}
The endpoint differences telescope. Moreover
\[
 \sum r=-4iS_Z,\qquad
 s+2i\cross{Z}{W}(Z+W)=W-Z,
 \qquad\sum s=-2i\mathcal T_Z=-4ivS_Z.
\]
Consequently the two linear terms together with the cubic term give
\begin{equation}\label{moment}
 \mathcal T(V)=-2S_Z\{t(Dv+3e/2+v|t|^2)
                +\overline t(D/2+H+ev/2)\}.
\end{equation}
These identities involve no division by $S(M)$.
For an arbitrary ordered polygon, translation by $c$ changes
$\mathcal T$ by $6cS$. This follows by expanding
$\cross{V_i+c}{V_{i+1}+c}$ and telescoping the coordinate-square and
coordinate-product endpoint differences.
Divide \eqref{moment} by $6S(M)$, translate by $z/2$ and use
\eqref{pedalarea} to obtain
\[
 C=\frac z2-
 \frac{z(D+2H+ev)+\overline z(2Dv+3e+v|z|^2/2)}
 {3(2D+|z|^2)}.
\]
Multiplying numerator and denominator by $d$ proves \eqref{centroid}.

\section{Definedness and an exact star example}
Positivity of $S(0)$ and $F$ makes \eqref{area} an exact classification of
the zero-area locus. For a simple winding-one family with $N\geq6$, the
first $N/2$ normals occupy a half-turn in circular order. Their squares are
at least three distinct circularly ordered unit-circle points. The full
$Z$-polygon traverses their convex polygon twice, so $S_Z>0$.
Equation \eqref{origin} implies $D>0$, hence $d>0$. Period four was
treated above. This completes Theorem \ref{main}.

The negative-$d$ case occurs in an exact primitive star. Take
\[
 a^2=221/96,\qquad b^2=27625/20736,\qquad\lambda=125/96.
\]
The caustic semiaxes are $1$ and $25/144$. Multiply the following normalized
pairs coordinatewise by $(a,b)$, retaining their order:
\begin{align*}
 &(0,1),\ (-12/13,-5/13),\ (1,0),\ (-12/13,5/13),\\
 &(0,-1),\ (12/13,5/13),\ (-1,0),\ (12/13,-5/13).
\end{align*}
The verification package checks exactly that every chord is tangent to the
caustic and obeys the specular reflection law. The eight vertices are
distinct, and an oriented-ray crossing count gives winding three.
This example was previously used in the repository's problem 5100065; no
novelty is claimed for its construction. Here
\[
 d=-120/221,\qquad F=2795/1728,\qquad v=7/13,\qquad
 -2F/d=123539/20736.
\]
For $M=(X,Y)$, direct Cartesian projection and signed shoelace give
\begin{align*}
 S(M)&=-\frac{\sqrt{30}}{134784}
          (20736(X^2+Y^2)-123539),\\
 \mathcal T_x(M)&=-\frac{\sqrt{30}X}{109512}
          (41472(X^2+Y^2)-272935),\\
 \mathcal T_y(M)&=-\frac{\sqrt{30}Y}{97344}
          (52992(X^2+Y^2)-347633).
\end{align*}
At $M=(13\sqrt{731}/144,0)$, $S=0$ but
$\mathcal T=(221\sqrt{21930}/10368,0)\ne0$.
Thus even a removable-zero interpretation cannot supply a centroid at
this diagnostic point.

\section*{Verification and disclosure}
The accompanying source and verification archive provides exact symbolic
checks of the displayed edge identities, unrestricted Cartesian star
polynomials and a separately derived real-coordinate certificate. Independent
orbit generators supply high-precision numerical controls, including stars,
large periods and near-degenerate caustics. Initial floating-point failures
and their precision repairs are recorded. Numerical tests corroborate, but
do not replace, the all-period proof above.

AI tools were used extensively in solving, verifying, literature research,
writing and adversarial review. Independent AI reviews are not external
human peer review or proof-assistant formalization. This is an unrefereed
research note; it has not received external human peer review. The bounded
priority reports, assumptions and reproducibility instructions accompany
the note. No claim of exclusive historical priority is made.

\begin{thebibliography}{9}
\bibitem{RGK} D. Reznik, R. Garcia and J. Koiller,
\emph{Fifty new invariants of N-periodics in the elliptic billiard},
Arnold Math. J. \textbf{7} (2021), 341--355.
\href{https://doi.org/10.1007/s40598-021-00174-y}{doi:10.1007/s40598-021-00174-y}.
See also \emph{Eighty new invariants},
\href{https://arxiv.org/abs/2004.12497v11}{arXiv:2004.12497v11}, Table 4.
\bibitem{BT} M. Bialy and S. Tabachnikov,
\emph{Dan Reznik's identities and more},
Eur. J. Math. \textbf{8} (2022), 1341--1354.
\href{https://doi.org/10.1007/s40879-020-00428-7}{doi:10.1007/s40879-020-00428-7}.
\bibitem{ST} R. E. Schwartz and S. Tabachnikov,
\emph{Centers of mass of Poncelet polygons, 200 years after},
Math. Intelligencer \textbf{38} (2016), no. 2, 29--34.
\href{https://arxiv.org/abs/1607.04766}{arXiv:1607.04766}.
\end{thebibliography}
\end{document}
